% =====================================================================
\chapter{Regla de la cadena y jacobiano}
\label{tema:2}
% =====================================================================

Una red neuronal es una composición de muchas funciones sencillas, y su pérdida depende de cada parámetro a través de largas cadenas de cálculos intermedios. Este tema generaliza la derivada a funciones con valores vectoriales mediante el jacobiano, demuestra la regla de la cadena en forma matricial y muestra cómo se organiza el cálculo sobre un grafo computacional, en modo directo o en modo inverso. Se sigue~\cite[§5.2, §5.3 y §5.6]{mml}, la clase~10 de~\cite{mit1802} y~\cite[§6.5]{gbc}.

\section{Funciones vectoriales y jacobiano}

\begin{definicion}[Jacobiano]\label{def:jacobiano}
Sea $U\subseteq\R^{n}$ abierto y $g=(g_1,\dots,g_m)\T\colon U\to\R^{m}$. La función $g$ es \emph{diferenciable} en $x$ si existe una matriz $J\in\R^{m\times n}$ tal que
\[
g(x+h)=g(x)+J\,h+o(\norm{h})\qquad\text{cuando } h\to 0.
\]
La matriz $J$ se llama \emph{jacobiano} de $g$ en $x$ y se escribe $J_g(x)$. Aplicando la \cref{prop:grad-parciales} a cada componente, $g$ es diferenciable si y solo si lo es cada $g_j$, y
\[
\bigl(J_g(x)\bigr)_{ji}=\frac{\partial g_j}{\partial x_i}(x),\qquad j=1,\dots,m,\quad i=1,\dots,n.
\]
La fila $j$ de $J_g(x)$ es $\nabla g_j(x)\T$; la columna $i$ es la derivada de $g$ respecto de $x_i$.
\end{definicion}

El jacobiano tiene una fila por cada salida y una columna por cada entrada. Para $m=1$ es una matriz fila y $\nabla f(x)=J_f(x)\T$.

\begin{ejemplo}\label{ej:jac-lineal}
Si $g(x)=Ax+b$ con $A\in\R^{m\times n}$ y $b\in\R^{m}$, entonces $g(x+h)=g(x)+Ah$ exactamente, luego $J_g(x)=A$ en todo punto. El jacobiano de una aplicación afín es su matriz.
\end{ejemplo}

\begin{ejemplo}\label{ej:jac}
Sea $g(x,y)=(xy,\;x^{2}+y)\T$. Entonces
\[
J_g(x,y)=\begin{pmatrix} y & x\\ 2x & 1\end{pmatrix},\qquad J_g(2,1)=\begin{pmatrix} 1 & 2\\ 4 & 1\end{pmatrix}.
\]
\end{ejemplo}

\section{La regla de la cadena}

\begin{teorema}[Regla de la cadena]\label{teo:cadena}
Sean $g\colon U\subseteq\R^{n}\to\R^{m}$ diferenciable en $x$ y $f\colon V\subseteq\R^{m}\to\R^{k}$ diferenciable en $g(x)$, con $g(U)\subseteq V$. Entonces $f\circ g$ es diferenciable en $x$ y
\[
J_{f\circ g}(x)=J_f\bigl(g(x)\bigr)\,J_g(x).
\]
\end{teorema}

\begin{proof}
Sean $A=J_g(x)$, $B=J_f(g(x))$ y $\Delta(h)=g(x+h)-g(x)=Ah+r_1(h)$, con $r_1(h)=o(\norm{h})$. Entonces $\norm{\Delta(h)}\le\norm{A}\,\norm{h}+\norm{r_1(h)}=O(\norm{h})$, donde $\norm{A}$ es la norma de operador. Por la diferenciabilidad de $f$ en $g(x)$,
\[
f\bigl(g(x+h)\bigr)-f\bigl(g(x)\bigr)=B\,\Delta(h)+r_2\bigl(\Delta(h)\bigr)=BA\,h+B\,r_1(h)+r_2\bigl(\Delta(h)\bigr),
\]
con $r_2(v)=o(\norm{v})$. El término $B\,r_1(h)$ es $o(\norm{h})$ porque $\norm{B\,r_1(h)}\le\norm{B}\,\norm{r_1(h)}$. Para el último, si $\Delta(h)\ne 0$ se escribe $\norm{r_2(\Delta)}/\norm{h}=\bigl(\norm{r_2(\Delta)}/\norm{\Delta}\bigr)\bigl(\norm{\Delta}/\norm{h}\bigr)$: el primer factor tiende a~0 porque $\Delta(h)\to 0$, y el segundo está acotado; si $\Delta(h)=0$, el término es nulo. Luego $f\circ g$ es diferenciable en $x$ con jacobiano $BA$.
\end{proof}

\begin{corolario}[Caso escalar]\label{cor:cadena-escalar}
Si $f\colon\R^{m}\to\R$ y $u=g(x)$, entonces
\[
\frac{\partial (f\circ g)}{\partial x_i}(x)=\sum_{j=1}^{m}\frac{\partial f}{\partial u_j}(u)\,\frac{\partial u_j}{\partial x_i}(x),
\qquad\text{es decir,}\qquad
\nabla_x (f\circ g)=J_g(x)\T\,\nabla_u f(u).
\]
\end{corolario}

\begin{proof}
Con $k=1$, el \cref{teo:cadena} da $J_{f\circ g}=\nabla f(u)\T J_g(x)$, una matriz fila cuya entrada $i$ es la suma del enunciado. Trasponiendo se obtiene la forma vectorial.
\end{proof}

La suma del \cref{cor:cadena-escalar} tiene una lectura directa: si $x_i$ influye en $f$ a través de varias variables intermedias $u_j$, su efecto total es la suma de los efectos transmitidos por cada una de ellas.

\begin{ejemplo}\label{ej:cadena}
Sea $f(u,v)=u^{2}v$ y $g$ la del \cref{ej:jac}, en el punto $(x,y)=(2,1)$.
\begin{enumerate}[label=\arabic*.]
  \item Valores intermedios: $u=xy=2$, $v=x^{2}+y=5$, $f=u^{2}v=20$.
  \item Gradiente exterior: $\partial f/\partial u=2uv=20$ y $\partial f/\partial v=u^{2}=4$.
  \item Por el \cref{cor:cadena-escalar},
  \[
  \nabla_{(x,y)}f=J_g(2,1)\T\begin{pmatrix}20\\4\end{pmatrix}=\begin{pmatrix}1&4\\2&1\end{pmatrix}\begin{pmatrix}20\\4\end{pmatrix}=\begin{pmatrix}36\\44\end{pmatrix}.
  \]
\end{enumerate}
Como comprobación, $f(g(x,y))=x^{2}y^{2}(x^{2}+y)=x^{4}y^{2}+x^{2}y^{3}$, cuyas parciales en $(2,1)$ son $4x^{3}y^{2}+2xy^{3}=32+4=36$ y $2x^{4}y+3x^{2}y^{2}=32+12=44$.
\end{ejemplo}

\section{Grafos computacionales}

\begin{definicion}[Grafo computacional]\label{def:grafo}
Un \emph{grafo computacional} es un grafo dirigido acíclico cuyos nodos son variables: los nodos sin predecesores son las entradas, y cada uno de los demás nodos $v$ se calcula como una función diferenciable $v=\phi_v(\text{padres}(v))$ de sus predecesores. Si la salida es un escalar $f$, se llama \emph{adjunto} de un nodo $v$ a $\bar v=\partial f/\partial v$, la derivada de $f$ respecto de $v$ considerando todas las variables posteriores como funciones de $v$~\cite[§6.5.1]{gbc}.
\end{definicion}

\begin{proposicion}[Retropropagación]\label{prop:retro}
En un grafo computacional con salida escalar $f$ se cumple $\bar f=1$ y, para todo nodo $v\neq f$,
\[
\bar v=\sum_{w\,:\,v\in\text{padres}(w)}\bar w\;\frac{\partial\phi_w}{\partial v}.
\]
Por tanto, recorriendo los nodos en orden topológico inverso se obtienen los adjuntos de todas las entradas, es decir, el gradiente de $f$.
\end{proposicion}

\begin{proof}
La variable $v$ influye en $f$ solo a través de los nodos $w$ de los que es padre. Considerando $f$ como función de esos nodos y estos como funciones de $v$, el \cref{cor:cadena-escalar} da la suma del enunciado. En orden topológico inverso, cuando se procesa $v$ ya se conocen todos los $\bar w$ de sus hijos.
\end{proof}

\begin{figure}[h]
\centering
\begin{tikzpicture}[>={Stealth[length=2mm]},font=\small,
  nodo/.style={draw,thick,rounded corners=2pt,minimum width=17mm,minimum height=7mm,align=center},
  adj/.style={->,dashed,gray!70!black}]
  \node[nodo] (x) at (0,1.3) {$x=2$};
  \node[nodo] (y) at (0,-1.3) {$y=1$};
  \node[nodo] (u) at (4,1.3) {$u=xy=2$};
  \node[nodo] (v) at (4,-1.3) {$v=x^{2}+y=5$};
  \node[nodo] (f) at (8,0) {$f=u^{2}v=20$};
  \draw[->,thick] (x) -- (u);
  \draw[->,thick] (y) -- (u);
  \draw[->,thick] (x) -- (v);
  \draw[->,thick] (y) -- (v);
  \draw[->,thick] (u) -- (f);
  \draw[->,thick] (v) -- (f);
  \node[above=1pt of u,gray!30!black] {$\bar u=2uv=20$};
  \node[below=1pt of v,gray!30!black] {$\bar v=u^{2}=4$};
  \node[above=1pt of f,gray!30!black] {$\bar f=1$};
  \node[above=1pt of x,gray!30!black] {$\bar x=\bar u\,y+\bar v\,2x=20+16=36$};
  \node[below=1pt of y,gray!30!black] {$\bar y=\bar u\,x+\bar v\cdot 1=40+4=44$};
\end{tikzpicture}
\caption{Grafo computacional del \cref{ej:cadena}. Dentro de cada nodo figura el valor calculado en el paso hacia delante; junto a él, en gris, el adjunto calculado en el paso hacia atrás con la \cref{prop:retro}. La entrada $x$ recibe contribuciones por dos caminos, que se suman.}
\label{fig:grafo}
\end{figure}

\section{Modo directo y modo inverso}

Sea una composición $F=f_L\circ\dots\circ f_1$, con $f_\ell\colon\R^{d_{\ell-1}}\to\R^{d_\ell}$. Por el \cref{teo:cadena},
\[
J_F=J_L\,J_{L-1}\cdots J_1,\qquad J_\ell\in\R^{d_\ell\times d_{\ell-1}},
\]
con cada $J_\ell$ evaluado en el valor intermedio correspondiente. El producto de matrices es asociativo, pero el coste de calcularlo depende del orden en que se agrupen los factores~\cite[§5.6]{mml}.

\begin{proposicion}[Coste de los dos modos]\label{prop:modos}
Multiplicar una matriz $a\times b$ por otra $b\times c$ requiere $abc$ productos escalares. En consecuencia:
\begin{enumerate}
  \item en \emph{modo directo}, se calcula $J_F v=J_L(\cdots(J_2(J_1 v)))$ para un vector $v\in\R^{d_0}$ con $\sum_\ell d_\ell d_{\ell-1}$ productos; obtener el jacobiano completo exige repetirlo $d_0$ veces, una por columna;
  \item en \emph{modo inverso}, se calcula $w\T J_F=((w\T J_L)J_{L-1})\cdots J_1$ para un vector $w\in\R^{d_L}$ con el mismo número de productos; obtener el jacobiano completo exige repetirlo $d_L$ veces, una por fila.
\end{enumerate}
Si la salida es escalar ($d_L=1$), el gradiente se obtiene con una sola pasada en modo inverso.
\end{proposicion}

\begin{proof}
Cada factor $J_\ell$ aplicado a un vector (por la derecha o por la izquierda) cuesta $d_\ell d_{\ell-1}$ productos, y en cada modo se realizan $L$ de estas operaciones. Con $v=e_i$ se obtiene la columna $i$ de $J_F$, y con $w=e_j$, la fila $j$. Con $d_L=1$, la única fila es $\nabla F\T$.
\end{proof}

La retropropagación de la \cref{prop:retro} es el modo inverso aplicado a un grafo. Su coste es un múltiplo pequeño del coste de evaluar $f$, independiente del número de entradas; por ello la pérdida de un modelo con miles de millones de parámetros puede derivarse con una sola pasada hacia atrás. El precio es la memoria: el paso hacia atrás necesita los valores intermedios del paso hacia delante. La implementación de este mecanismo (diferenciación automática) es el objeto del módulo~10.

\begin{observacion}[Conexión con los LLM]
Un Transformer es un grafo computacional con miles de millones de nodos. Su salida escalar es la pérdida media del lote, y las bibliotecas como PyTorch o JAX registran el grafo durante el paso hacia delante para recorrerlo después en orden inverso. Ninguna de ellas forma explícitamente un jacobiano: solo calculan productos de la forma $w\T J_\ell$.
\end{observacion}

\begin{lecturas}
\item Deisenroth, Faisal y Ong~\cite{mml}, §5.2 (regla de la cadena), §5.3 «Gradients of Vector-Valued Functions» y §5.6 «Backpropagation and Automatic Differentiation».
\item MIT 18.02~\cite{mit1802}, clase~10 «Differentials; Chain Rule».
\item Goodfellow, Bengio y Courville~\cite{gbc}, §6.5.1 «Computational Graphs», §6.5.2 «Chain Rule of Calculus» y §6.5.3 «Recursively Applying the Chain Rule to Obtain Backprop».
\item Zhang \emph{et al.}~\cite{d2l}, §2.4.4 «Chain Rule».
\end{lecturas}

\section*{Ejercicios del tema 2}
\addcontentsline{toc}{section}{Ejercicios del tema 2}

\begin{ejercicio}{1}{eje:2-formas}
Sean $g\colon\R^{5}\to\R^{3}$, $h\colon\R^{3}\to\R^{2}$ y $f\colon\R^{2}\to\R$ diferenciables. Indique la forma de $J_g$, $J_h$, $J_f$ y $J_{f\circ h\circ g}$, y la del gradiente de $f\circ h\circ g$. Compruebe que las dimensiones del producto del \cref{teo:cadena} encajan.
\end{ejercicio}

\begin{ejercicio}{2}{eje:2-producto}
Sea $F(x,y)=(x+y^{2})(x^{2}-y)$. Calcule $\partial F/\partial x$ y $\partial F/\partial y$ en $(1,2)$ con la regla de la cadena, tomando $u=x+y^{2}$ y $v=x^{2}-y$, y compruebe el resultado desarrollando $F$.
\end{ejercicio}

\begin{ejercicio}{2}{eje:2-jacobiano}
Sea $g(x,y)=(x^{2}y,\;3x-y,\;xy)\T$. Escriba $J_g(1,2)$ e indique su forma. Use $J_g\T$ para obtener el gradiente de $s=g_1+g_2+g_3$ en $(1,2)$.
\end{ejercicio}

\begin{ejercicio}{2}{eje:2-d2l}
Sea $w=f(x,y,z)$, donde $x$, $y$ y $z$ son funciones de dos variables $a$ y $b$. (a) Escriba la regla de la cadena para $\partial w/\partial a$ y $\partial w/\partial b$. (b) Aplíquela con $f(x,y,z)=xyz$, $x=a+b$, $y=a-b$, $z=ab$, en $(a,b)=(2,1)$, y compruebe el resultado expresando $w$ directamente en función de $a$ y $b$.
\fuente{adaptado de Zhang \emph{et al.}~\cite{d2l}, §2.4, ejercicio~9 (se añade el caso numérico).}
\end{ejercicio}

\begin{ejercicio}{2}{eje:2-coste}
Sea $F=f_3\circ f_2\circ f_1$ con $f_1\colon\R^{n}\to\R^{d}$, $f_2\colon\R^{d}\to\R^{d}$ y $f_3\colon\R^{d}\to\R$. Cuente los productos escalares necesarios para obtener $\nabla F$ multiplicando los jacobianos (a) de derecha a izquierda, $J_3(J_2J_1)$, y (b) de izquierda a derecha, $(J_3J_2)J_1$. Evalúe ambos con $n=d=1000$.
\end{ejercicio}

\begin{ejercicio}{3}{eje:2-escalar}
Sean $a,b\colon\R^{n}\to\R^{m}$ diferenciables y $s(x)=a(x)\T b(x)$. Demuestre, con el \cref{teo:cadena}, que $\nabla s(x)=J_a(x)\T b(x)+J_b(x)\T a(x)$. Deduzca el gradiente de la forma cuadrática $q(x)=x\T Ax$ para $A\in\R^{n\times n}$ y particularícelo al caso simétrico.
\end{ejercicio}
